% GENERATED FILE. Edit canonical lessons, then run npm run generate:books.
% canonical-source-hash: fnv1a64:e6e8fb9681acffa7
% canonical-lessons: LA-C318-cremo, LA-C318-uro, LA-C318-flagro, LA-C318-fumo, LA-C318-conflagro

\chapter{Ārdet}
\label{ch:la-ardet}
\hlchaptermodality{318}

\begin{chapteropening}
\textbf{By the end of this chapter:} \emph{I can say how things burned: to burn up, to scorch, to blaze, to smoke and to go up in flames.}

Complete the last lesson of chapter 318: I can say how things burned: to burn up, to scorch, to blaze, to smoke and to go up in flames.
\end{chapteropening}

\section[cremō, cremāre]{cremō, cremāre — to burn up, to cremate}

\label{lesson:LA-C318-cremo}



\begin{quote}
Before the new one: say the Latin for a firebrand, then the Latin for burning heat.
\end{quote}

\subsection*{You'll want to know: cremō, cremāre}
\textbf{cremō, cremāre} — \textquotedblleft{}to burn up, to cremate\textquotedblright{}.

\begin{quote}
Ignis tōtum vīcum cremāvit.
\end{quote}

\begin{quote}
\emph{The fire burned up the whole street.}
\end{quote}

\emph{cremō, cremāre, cremāvī, cremātum}. Romans usually burned their dead, and English \emph{cremate} comes from this verb.

\begin{cousinweb}
\textbf{cremō} is inherited; EtymDB traces it to a Proto-Indo-European root meaning \textquotedblleft{}to burn\textquotedblright{}, also seen in \emph{carbō}, \textquotedblleft{}coal\textquotedblright{}. English \emph{cremate} comes from it.
\end{cousinweb}

\begin{grammarlens}[title={What you've built}]
One more word for telling, explaining and describing at B1.
\end{grammarlens}

\subsection*{Your turn}
\begin{itemize}
\raggedright
  \item \emph{Say it:} \emph{cremō, cremāre}
  \item \emph{Say it:} \emph{cremō, cremāre}, once more
  \item \emph{Say it:} a sentence of your own with \emph{cremō, cremāre}
  \item \emph{Recall:} say the Latin for a firebrand, then the Latin for burning heat, then say \emph{cremō, cremāre} again
\end{itemize}

\begin{tcolorbox}[breakable,colback=teal!4,colframe=teal!35!black,title={Before you move on}]
What does cremō, cremāre mean? (\textquotedblleft{}To burn up, to cremate\textquotedblright{}.) Say it once more.
\end{tcolorbox}
\section[ūrō, ūrere]{ūrō, ūrere — to burn; to scorch}

\label{lesson:LA-C318-uro}



\begin{quote}
Before the new one: say the Latin for burning heat, then the Latin for to burn up.
\end{quote}

\subsection*{You'll want to know: ūrō, ūrere}
\textbf{ūrō, ūrere} — \textquotedblleft{}to burn; to scorch\textquotedblright{}.

\begin{quote}
Sōl agrōs ūrēbat.
\end{quote}

\begin{quote}
\emph{The sun was scorching the fields.}
\end{quote}

\emph{ūrō, ūrere, ussī, ustum}. It has many compounds; English \emph{combustion} comes from \emph{combūrō}.

\begin{cousinweb}
\textbf{ūrō} is an inherited verb; Lewis \& Short compare Greek \emph{heuō}, \textquotedblleft{}to singe\textquotedblright{}. English \emph{combustion} comes from its compound \emph{combūrō}.
\end{cousinweb}

\begin{grammarlens}[title={What you've built}]
One more word for telling, explaining and describing at B1.
\end{grammarlens}

\subsection*{Your turn}
\begin{itemize}
\raggedright
  \item \emph{Say it:} \emph{ūrō, ūrere}
  \item \emph{Say it:} \emph{ūrō, ūrere}, once more
  \item \emph{Say it:} a sentence of your own with \emph{ūrō, ūrere}
  \item \emph{Recall:} say the Latin for burning heat, then the Latin for to burn up, then say \emph{ūrō, ūrere} again
\end{itemize}

\begin{tcolorbox}[breakable,colback=teal!4,colframe=teal!35!black,title={Before you move on}]
What does ūrō, ūrere mean? (\textquotedblleft{}To burn; to scorch\textquotedblright{}.) Say it once more.
\end{tcolorbox}
\section[flagrō, flagrāre]{flagrō, flagrāre — to blaze, to be on fire}

\label{lesson:LA-C318-flagro}



\begin{quote}
Before the new one: say the Latin for to burn up, then the Latin for to burn.
\end{quote}

\subsection*{You'll want to know: flagrō, flagrāre}
\textbf{flagrō, flagrāre} — \textquotedblleft{}to blaze, to be on fire\textquotedblright{}.

\begin{quote}
Domus flagrābat, et vīcīnī aquam portābant.
\end{quote}

\begin{quote}
\emph{The house was blazing, and the neighbours were carrying water.}
\end{quote}

\emph{flagrō, flagrāre, flagrāvī}: to blaze, and of people, to burn with a feeling. English \emph{flagrant}.

\begin{cousinweb}
\textbf{flagrō} is inherited from Proto-Italic \emph{*flagros}, \textquotedblleft{}burning\textquotedblright{} (EtymDB), of the same family as \emph{flamma}, \textquotedblleft{}flame\textquotedblright{}. English \emph{flagrant} comes from it.
\end{cousinweb}

\begin{grammarlens}[title={What you've built}]
One more word for telling, explaining and describing at B1.
\end{grammarlens}

\subsection*{Your turn}
\begin{itemize}
\raggedright
  \item \emph{Say it:} \emph{flagrō, flagrāre}
  \item \emph{Say it:} \emph{flagrō, flagrāre}, once more
  \item \emph{Say it:} a sentence of your own with \emph{flagrō, flagrāre}
  \item \emph{Recall:} say the Latin for to burn up, then the Latin for to burn, then say \emph{flagrō, flagrāre} again
\end{itemize}

\begin{tcolorbox}[breakable,colback=teal!4,colframe=teal!35!black,title={Before you move on}]
What does flagrō, flagrāre mean? (\textquotedblleft{}To blaze, to be on fire\textquotedblright{}.) Say it once more.
\end{tcolorbox}
\section[fūmō, fūmāre]{fūmō, fūmāre — to smoke}

\label{lesson:LA-C318-fumo}



\begin{quote}
Before the new one: say the Latin for to burn, then the Latin for to blaze.
\end{quote}

\subsection*{You'll want to know: fūmō, fūmāre}
\textbf{fūmō, fūmāre} — \textquotedblleft{}to smoke\textquotedblright{}.

\begin{quote}
Post incendium ruīnae diū fūmābant.
\end{quote}

\begin{quote}
\emph{After the fire the ruins smoked for a long time.}
\end{quote}

\emph{fūmō, fūmāre, fūmāvī}, from \emph{fūmus}, 'smoke' (taught). English \emph{fume}.

\begin{cousinweb}
\textbf{fūmō} is built on \emph{fūmus}, \textquotedblleft{}smoke\textquotedblright{} (EtymDB). English \emph{fume} comes from the same family through French.
\end{cousinweb}

\begin{grammarlens}[title={What you've built}]
One more word for telling, explaining and describing at B1.
\end{grammarlens}

\subsection*{Your turn}
\begin{itemize}
\raggedright
  \item \emph{Say it:} \emph{fūmō, fūmāre}
  \item \emph{Say it:} \emph{fūmō, fūmāre}, once more
  \item \emph{Say it:} a sentence of your own with \emph{fūmō, fūmāre}
  \item \emph{Recall:} say the Latin for to burn, then the Latin for to blaze, then say \emph{fūmō, fūmāre} again
\end{itemize}

\begin{tcolorbox}[breakable,colback=teal!4,colframe=teal!35!black,title={Before you move on}]
What does fūmō, fūmāre mean? (\textquotedblleft{}To smoke\textquotedblright{}.) Say it once more.
\end{tcolorbox}
\section[cōnflagrō, cōnflagrāre]{cōnflagrō, cōnflagrāre — to go up in flames}

\label{lesson:LA-C318-conflagro}



\begin{quote}
Before the new one: say the Latin for to blaze, then the Latin for to smoke.
\end{quote}

\subsection*{You'll want to know: cōnflagrō, cōnflagrāre}
\textbf{cōnflagrō, cōnflagrāre} — \textquotedblleft{}to go up in flames\textquotedblright{}.

\begin{quote}
Tōtum aedificium ūnā nocte cōnflagrāvit.
\end{quote}

\begin{quote}
\emph{The whole building went up in flames in a single night.}
\end{quote}

\emph{cōnflagrō, cōnflagrāre, cōnflagrāvī, cōnflagrātum}: \emph{con-} 'wholly' + \emph{flagrō}. English \emph{conflagration}.

\begin{cousinweb}
\textbf{cōnflagrō} is built on \emph{flagrō}, \textquotedblleft{}to blaze\textquotedblright{}, with \emph{con-}, \textquotedblleft{}wholly\textquotedblright{} (Lewis \& Short): to burn up entirely. English \emph{conflagration} comes from it.
\end{cousinweb}

\begin{grammarlens}[title={What you've built}]
One more word for telling, explaining and describing at B1.
\end{grammarlens}

\subsection*{Your turn}
\begin{itemize}
\raggedright
  \item \emph{Say it:} \emph{cōnflagrō, cōnflagrāre}
  \item \emph{Say it:} \emph{cōnflagrō, cōnflagrāre}, once more
  \item \emph{Say it:} a sentence of your own with \emph{cōnflagrō, cōnflagrāre}
  \item \emph{Recall:} say the Latin for to blaze, then the Latin for to smoke, then say \emph{cōnflagrō, cōnflagrāre} again
\end{itemize}

\begin{tcolorbox}[breakable,colback=teal!4,colframe=teal!35!black,title={Before you move on}]
What does cōnflagrō, cōnflagrāre mean? (\textquotedblleft{}To go up in flames\textquotedblright{}.) Say it once more.
\end{tcolorbox}
